\documentclass[10.5pt]{article}
\usepackage[letterpaper, margin=0.6in, top=0.5in, bottom=0.5in]{geometry}
\usepackage{amsmath,amssymb,amsthm}
\usepackage{fancyhdr}
\usepackage{xcolor}
\usepackage{tikz}
\usepackage{enumitem}
\usepackage{tcolorbox}
\usepackage{hyperref}

% Define colors matching instructor board annotations
\definecolor{chalkblue}{RGB}{0, 105, 185}
\definecolor{chalkred}{RGB}{205, 35, 35}
\definecolor{chalkpink}{RGB}{230, 90, 120}
\definecolor{chalkcyan}{RGB}{15, 160, 190}
\definecolor{chalkgreen}{RGB}{20, 130, 50}
\definecolor{boxbg}{RGB}{252, 252, 254}
\definecolor{boxframe}{RGB}{150, 160, 180}

\pagestyle{fancy}
\fancyhf{}
\renewcommand{\headrulewidth}{0pt}
\renewcommand{\footrulewidth}{0pt}

\setlength{\parindent}{0pt}
\setlength{\parskip}{2.5pt}

\begin{document}

% ==============================================================================
% PAGE 1
% ==============================================================================
\noindent \textbf{\large Math 105L} \hfill \textbf{\large Lecture 2-1}\\[1pt]
{\textbf{\Large Power Functions, Monomial Ordering, \& Polynomials}}\\[-5pt]
\rule{\textwidth}{1.2pt}

\vspace{0.03in}

\textbf{\large 1. Power Functions}

A \textbf{power function} is any function of the form:
\[
  \mathbf{f(x) = k x^p}
\]
where the coefficient $k$ and exponent $p$ are real constants.
\begin{itemize}[leftmargin=1.2em, topsep=1pt, itemsep=2pt]
  \item \textbf{Radicals as Fractional Powers:} $\sqrt{x} = \mathbf{x^{1/2}}, \quad \sqrt[3]{x} = \mathbf{x^{1/3}}$.
  \item \textbf{Reciprocals as Negative Powers:} $\frac{1}{\sqrt[3]{x^7}} = \frac{1}{x^{7/3}} = \mathbf{x^{-7/3}}$.
\end{itemize}

\vspace{0.03in}

\begin{tcolorbox}[colback=boxbg, colframe=boxframe, boxrule=0.8pt, arc=3pt, left=8pt, right=8pt, top=4pt, bottom=4pt]
\textbf{\large Relative Ordering \& Geometry of Monomials $x^n$ ($n \ge 1$)}
All integer monomial curves $y = x^n$ pass through the common points $\mathbf{(0,0)}$ and $\mathbf{(1,1)}$ in the first quadrant:
\begin{itemize}[leftmargin=1.2em, topsep=2pt, itemsep=4pt]
  \item \textbf{Inside the Unit Interval $(0, 1)$:}
  Smaller powers are \textbf{strictly greater} than higher powers because multiplying numbers between $0$ and $1$ reduces their size:
  \[
    \mathbf{x > x^2 > x^3 > x^4 > x^5 \quad \text{for } x \in (0, 1)}
  \]
  \item \textbf{Beyond the Unit Interval $(1, \infty)$:}
  Higher powers grow much faster and \textbf{dominate}:
  \[
    \mathbf{x^5 > x^4 > x^3 > x^2 > x \quad \text{for } x \in (1, \infty)}
  \]
  \item \textbf{Symmetry Across the Axes:}
  \begin{itemize}[label=--]
    \item \textbf{Even Powers ($x^2, x^4$):} Even functions ($f(-x) = f(x)$), symmetric across the $y$-axis. Both arms point up ($+\infty$ as $x \to \pm\infty$).
    \item \textbf{Odd Powers ($x, x^3, x^5$):} Odd functions ($f(-x) = -f(x)$), symmetric through the origin. Passes through $(-1, -1)$.
  \end{itemize}
\end{itemize}
\end{tcolorbox}

\vspace{0.04in}

\textbf{\large 2. Polynomials \& Degree Classification}

\begin{tcolorbox}[colback=chalkblue!6, colframe=chalkblue, boxrule=1pt, arc=3pt, left=8pt, right=8pt, top=4pt, bottom=4pt]
\textbf{Definition of a Polynomial:}\\[1pt]
A \textbf{polynomial} is a sum of power functions with non-negative integer exponents:
\[
  \mathbf{p(x) = a_n x^n + a_{n-1} x^{n-1} + \dots + a_1 x + a_0 \quad (a_n \ne 0)}
\]
The largest exponent $n$ is the \textbf{degree} of the polynomial:
\begin{center}
\begin{tabular}{lcl|lcl}
  $n = 1$: & \textbf{Linear} & ($y = mx + b$) & $n = 2$: & \textbf{Quadratic} & ($ax^2 + bx + c$) \\
  $n = 3$: & \textbf{Cubic} & ($ax^3 + \dots$) & $n = 4$: & \textbf{Quartic} & ($ax^4 + \dots$) \\
  $n = 5$: & \textbf{Quintic} & ($ax^5 + \dots$) & $n \ge 6$: & \textbf{Degree-$n$} & \\
\end{tabular}
\end{center}
\end{tcolorbox}

\vspace{0.03in}

\textbf{End Behavior Dominance:}
As $x \to \pm\infty$, the highest degree term $a_n x^n$ vastly outgrows all lower terms combined, determining the ultimate end-behavior of the polynomial.

\end{document}

